% Generated draft: requires mathematical and source review.
\documentclass[11pt]{article}
\usepackage[margin=25mm]{geometry}
\usepackage{fontspec}
\setmainfont{Latin Modern Roman}
\usepackage{amsmath,amssymb,mathtools,unicode-math}
\setmathfont{Latin Modern Math}
\usepackage{enumitem,needspace,xurl,hyperref}
\setlength{\parindent}{0pt}
\setlength{\parskip}{5pt}
\setlength{\emergencystretch}{3em}
\newcommand{\N}{\mathbb{N}}
\newcommand{\Z}{\mathbb{Z}}
\newcommand{\Q}{\mathbb{Q}}
\newcommand{\R}{\mathbb{R}}
\newcommand{\C}{\mathbb{C}}
\newcommand{\F}{\mathbb{F}}
\newcommand{\A}{\mathbb{A}}
\newcommand{\PP}{\mathbb{P}}
\newcommand{\E}{\mathbb{E}}
\newcommand{\eps}{\varepsilon}
\newcommand{\dd}{\,\mathrm{d}}
\begin{document}
\section*{UnsolvedMath ID: 30002332}
\textbf{Source problem code:} catalogue entry
\subsection{Definitions and mathematical statement}
For a standard geometric interpretation, let a knot type mean an ambient isotopy class of tame embeddings of the circle in \(\R^3\). Define its crossing number \(\operatorname{Cr}(K)\) as the minimum number of crossings in a regular planar diagram representing \(K\). For an embedded closed \(C^{1,1}\) curve \(\gamma\), let \(L(\gamma)\) be its length and let its thickness \(\tau(\gamma)\) be its reach: the supremum of radii \(r>0\) for which every point at distance less than \(r\) from \(\gamma\) has a unique nearest point on \(\gamma\). Set
\[
\operatorname{Rop}(K)=\inf_{\gamma\text{ representing }K,\ \tau(\gamma)>0}\frac{L(\gamma)}{\tau(\gamma)}.
\]
The tame-knot domain, regularity, and thickness convention make the questions precise but are not explicitly specified in the supplied record. A knot is alternating if it admits a diagram in which overcrossings and undercrossings alternate while traversing the knot.

The main question, interpreting the asymptotic constant uniformly over knot types, is whether
\[
\exists C>0\ \exists n_0\in\N,\ n_0\ge1\ \forall K:\quad
\operatorname{Cr}(K)\ge n_0\ \Longrightarrow\ \operatorname{Rop}(K)\le C\operatorname{Cr}(K).
\]
The crossing threshold expresses an asymptotic assertion and avoids imposing a zero upper bound on the unknot, whose crossing number is zero.

If this bound fails, the supplied record asks whether infinitely many knots have faster growth. A precise candidate interpretation is the existence of pairwise distinct knot types \((K_j)_{j\ge1}\), with positive crossing numbers, such that
\[
\operatorname{Cr}(K_j)\longrightarrow\infty,
\qquad
\frac{\operatorname{Rop}(K_j)}{\operatorname{Cr}(K_j)}\longrightarrow\infty.
\]
The source does not specify a family, a growth rate, or whether this precise sequential interpretation is intended. Its original wording calls this question ``equivalent'' to the upper-bound question; the faster-growth alternative concerns failure of the upper bound, rather than its affirmative assertion.

The alternating-knot question, interpreted as a uniform lower bound over all alternating knot types, is whether
\[
\exists c>0\ \exists n_1\in\N,\ n_1\ge1\ \forall K\text{ alternating}:\quad
\operatorname{Cr}(K)\ge n_1\ \Longrightarrow\ \operatorname{Rop}(K)\ge c\operatorname{Cr}(K).
\]
No restriction to prime alternating knots is supplied. Here ``linear-order lower bound'' is interpreted as \(\Omega(\operatorname{Cr}(K))\), without asserting a matching upper bound. These formulations are statement-only interpretations requiring mathematical review.
\subsection{Short English statement}
Ask whether a universal constant bounds knot ropelength above by crossing number, whether failure yields a sequence with unbounded ropelength-to-crossing-number ratio, and whether alternating knots satisfy a uniform linear lower bound. The domain, geometric conventions, and meaning of faster growth require confirmation. The supplied open status is dated curation, not independently verified.
\subsection{Sources}
\begin{itemize}
\item Supplied local UnsolvedMath registry record, permanent ID \(30002332\), title \emph{Linear Ropelength Bounds in Crossing Number}; collection: Oberwolfach Reports Open Problems; exact source problem code: \texttt{catalogue entry}. The record gives proposed year 2013 and identifies no individual proposer.
\item Supplied source citation: \emph{Geometric Knot Theory (2014)}, Oberwolfach Reports, pp.~1313--1358, DOI \texttt{10.4171/owr/2013/22}. \url{https://doi.org/10.4171/owr/2013/22}. The background describes the reports as documenting mathematical workshops and their open-problem discussions. The citation's 2014 label and the record's 2013 report designation require reconciliation.
\item Supplied literature reference: Diao--Ernst--Por--Ziegler, \emph{The ropelength of a knot}. \url{https://arxiv.org/abs/0912.3282}.
\item Supplied literature reference: Oberwolfach Report 22/2013. \url{https://doi.org/10.4171/owr/2013/22}.
\item Literature assessment dated exactly \texttt{2026-08-21}: the supplied curation labels the problem open, reports that the best general theorem located gives \(O(\operatorname{Cr}(K)\log^5\operatorname{Cr}(K))\), and describes the linear upper bound and listed family questions as unresolved. This is supplied dated curation derived from the OWR material, not an independently verified literature assessment; its status and cited claims require independent checking.
\item The registry labels the statement \texttt{corrected\_verified} and reports a manual semantic, typographic, and LaTeX review of the cleanup. This supplied editorial label does not establish mathematical correctness or independent verification of the cited sources.
\end{itemize}
\end{document}
